\documentclass{beamer}
\usepackage[spanish]{babel}
\usepackage[utf8]{inputenc}
\usepackage{graphicx}
\usepackage{mathdots} % para el comando \iddots
\usepackage{amssymb, amsmath, amsbsy} % simbolitos
\usepackage{xcolor}


\begin{document}

\begin{frame}
    \frametitle{Ejemplos y definiciones de fen\'omenos din\'amicos}
    
    Un sistema dinámico puede ser cualquier mecanismo que evolucione de forma determinista en el tiempo.
Se pueden encontrar ejemplos simples en mecánica, uno puede pensar en el péndulo o el
sistema solar. Sistemas similares ocurren en química y meteorología. Debemos notar
que en los sistemas biológicos y económicos es menos claro cuando se trata de
determinismo.1 Nos interesan las \emph{evoluciones}, es decir, las funciones que describen el
\emph{estado} de un sistema en función del \emph{tiempo} y que satisfacen las \emph{ecuaciones de movimiento} del
sistema. Las definiciones matemáticas de estos conceptos siguen más adelante.

\end{frame}

\begin{frame}

The simplest type of evolution is \emph{stationary}, where the state is constant in time.
Next we also know \emph{periodic evolutions}. Here, after a fixed period, the system always
returns to the same state and we get a precise repetition of what happened in the
previous period. Stationary and periodic evolutions are very regular and predictable.
Here prediction simply is based on matching with past observations. Apart from
these types of evolutions, in quite simple systems one meets evolutions that are not
so regular and predictable. In cases where the unpredictability can be established in
a definite way, we speak of \emph{chaotic} behaviour.

\end{frame}

\begin{frame}

En la primera parte de este capítulo, utilizando sistemas como el péndulo con o sin
amortiguamiento o forzamiento externo, damos una impresión de los tipos de evolución
que puede ocurrir. Mostramos que por lo general un sistema dado tiene varios tipos de evolución.
Qué tipos son típicos o predominantes depende en gran medida del tipo de sistema
a mano; por ejemplo, en sistemas mecánicos es importante si consideramos
disipación de energía, por ejemplo, por fricción o amortiguamiento. Durante esta exposición
el marco matemático en el que describir los diversos fenómenos se convierte en
más claro Esto determina el lenguaje y la forma de pensar de la disciplina de
Sistemas Dinámicos.

\end{frame}

\begin{frame}

En la segunda parte de este capítulo somos más explícitos, dando una definición formal
del sistema dinámico. Entonces conceptos tales como estado, tiempo y similares también son
discutido Después de esto, volviendo a los ejemplos anteriores, ilustramos estos conceptos.
En la parte final del capítulo tratamos una serie de ejemplos misceláneos que regresan regularmente en partes posteriores del libro. Aquí prestamos especial atención a
ejemplos que, históricamente hablando, han orientado el desarrollo de la
disciplina. Por lo tanto, estos ejemplos actúan como chistes a lo largo del libro.


\end{frame}

\begin{frame}\frametitle{1.1 The pendulum as a dynamical system}

Como preparación para las definiciones formales que veremos más adelante en este capítulo, primero trataremos el péndulo en varias variaciones: con o sin amortiguamiento y forzamiento externo. Presentamos gráficos de evoluciones calculadas numéricamente. De estos gráficos queda claro que existen diferencias cualitativas entre las diversas evoluciones y que el tipo de evolución que es típica o predominante depende en gran medida del contexto en cuestión. Esto ya es válido para la clase restringida de sistemas mecánicos considerados
aquí.


\end{frame}

\begin{frame}\frametitle{1.1.1 El péndulo libre}

El péndulo matemático plano consta de una varilla, suspendida en un punto fijo en
un plano vertical en el que el péndulo puede moverse. Se piensa que toda masa es
concentrada en una masa puntual al final de la varilla (ver Figura 1.1), y la varilla misma
es sin masa. También se supone que la varilla está rígida. El péndulo tiene masa m y longitud `.
Suponemos además que la suspensión es tal que el péndulo no sólo puede
oscilar, pero también puede ir 'sobre la parte superior'. En el caso sin forzamiento externo,
hablar del péndulo libre.

\end{frame}\frametitle{1.1.1.1 El péndulo libre no amortiguado}

\begin{frame}

En el caso sin amortiguamiento y forzado el péndulo solo está sujeto a la gravedad,
con aceleración $g$: La fuerza gravitacional apunta verticalmente hacia abajo con
fuerza $mg$ y tiene un componente $-mg \sin \varphi$ a lo largo del círculo descrito por el
masa puntual; consulte la figura 1.1. Aquí $\varphi$ es el ángulo entre la barra y el hacia abajo
vertical, a menudo llamada "desviación", expresada en radianes. La distancia del punto
masa desde la 'posición de reposo' ($\varphi = 0$), medida a lo largo del círculo de todos sus posibles
posiciones, entonces es $l \varphi$. La relación entre fuerza y movimiento está determinada por
la famosa ley de Newton2 $F=ma$,  donde F denota la fuerza, m la masa y $a$ la
aceleración.



\end{frame}


\begin{frame}

Por la rigidez de la varilla, no se produce movimiento en la dirección radial y
por lo tanto, simplemente aplique la ley de Newton en la $\varphi$-dirección. La componente de la fuerza en
la $\varphi$-dirección viene dada por $-mg \sin \varphi$, donde el signo menos representa el hecho
que la fuerza hace retroceder el péndulo a $\varphi = 0$. Para la aceleración $a$ tenemos

\[
a = \frac{d^2 (l \varphi)}{d t^2} = l \frac{d^2 \varphi}{dt^2}
\]


donde ($d^2 \varphi / dt^2)(t) = \varphi''(t)$ es la segunda derivada de la función $ t \to \varphi(t)$.

\end{frame}


\begin{frame}

Sustituida en la ley de Newton esto da

\[
ml \varphi'' = -mg \sin \varphi
\]

o equivalente,

\[
\varphi'' = - \frac{g}{l} \sin \varphi,
\]

donde obviamente $m, l >0$. Así derivamos la ecuación de movimiento del péndulo,
que en este caso es una ecuación diferencial ordinaria. Esto significa que las evoluciones
están dadas por las funciones $t \to \varphi(t)$ que satisfacen la ecuación de movimiento (1.1). En el
continuación abreviamos $\omega = \sqrt{g/l}$.

\end{frame}


\begin{frame}


\end{frame}


\begin{frame}


\end{frame}


\begin{frame}


\end{frame}


\begin{frame}


\end{frame}


\begin{frame}


\end{frame}







\end{document}
